% ======================================================================
\chapter[Магнитный момент. Опыт Штерна--Герлаха. Спин]%
        {Магнитный момент электрона. Опыт Штерна--Герлаха. Спин}
\label{t:12}
% ======================================================================

\begin{qbox}
\textbf{Билет 12.} Магнитный орбитальный момент электронов, гиромагнитное отношение, $g$-фактор,
магнетон Бора. Опыт Штерна--Герлаха. Спин электрона, спиновый $g$-фактор.
\end{qbox}

\section{Содержание билета}

\textbf{Подтемы:}
\begin{enumerate}
 \item орбитальный магнитный момент: классический вывод, связь с механическим моментом;
 \item гиромагнитное отношение, $g$-фактор, магнетон Бора; квантование проекции магнитного момента;
 \item атом в магнитном поле: энергия, ларморова прецессия, эффект Зеемана (кратко);
 \item опыт Штерна--Герлаха: идея, расчёт отклонения, результат и его интерпретация;
 \item спин электрона: собственный момент, $s=\frac12$, спиновый магнитный момент, $g_s\approx2$; полный момент.
\end{enumerate}

\textbf{Необходимые понятия:} магнитный момент тока $\mu=IS/c$; энергия и сила для диполя в
поле; оператор момента импульса и его собственные значения (\tref{t:09}).

\section{Определения и физические принципы}

\begin{definition}[магнитный момент и гиромагнитное отношение]
Магнитный момент $\vect\mu$ — величина, определяющая энергию $U=-\vect\mu\cdot\vect B$ и момент сил
$\vect M=\vect\mu\times\vect B$, действующие на систему во внешнем магнитном поле. \emph{Гиромагнитное
отношение} $\gamma$ — коэффициент пропорциональности между магнитным и механическим моментами:
$\vect\mu=\gamma\vect L$.
\end{definition}

\begin{definition}[$g$-фактор и магнетон Бора]
Для электрона магнитный момент записывают в виде
\begin{equation}\label{eq:12:g}
  \vect\mu=-g\,\muB\,\frac{\vect J}{\hbar},\qquad
  \muB=\frac{e\hbar}{2\mel c}=9{,}274\cdot10^{-21}\ \frac{\text{эрг}}{\text{Гс}}=5{,}788\cdot10^{-5}\ \frac{\text{эВ}}{\text{Тл}},
\end{equation}
где $\vect J$ — соответствующий механический момент, $g$ — безразмерный $g$-фактор (фактор Ланде).
Для орбитального движения $g_l=1$, для спина $g_s\approx2$. \emph{Магнетон Бора} $\muB$ — естественная единица
магнитного момента электронной оболочки (в СИ $\muB=e\hbar/2\mel$).
\end{definition}

\begin{definition}[спин]
\emph{Спин} — собственный (внутренний) момент импульса частицы, не связанный с её движением в
пространстве. Оператор спина $\op{\vect S}$ подчиняется тем же коммутационным соотношениям, что и $\op{\vect L}$;
$S^2=\hbar^2s(s+1)$, $S_z=m_s\hbar$, $m_s=-s,\dots,s$. Для электрона $s=\frac12$: $|\vect S|=\frac{\sqrt3}{2}\hbar$, $S_z=\pm\frac\hbar2$.
\end{definition}

\emph{Физический смысл.} Магнитные свойства атомов определяются механическими моментами электронов.
Квантование проекции момента ведёт к квантованию проекции магнитного момента — «пространственному
квантованию», непосредственно наблюдаемому в опыте Штерна--Герлаха. Спин — релятивистское по
происхождению (следует из уравнения Дирака) квантовое свойство без классического аналога: модель
«вращающегося шарика» приводит к линейным скоростям поверхности, много большим $c$.

\section{Основная теория}

\subsection{Орбитальный магнитный момент}

Электрон, движущийся по круговой орбите, — ток $I=e/T=ev/2\pi r$. Его магнитный момент
$\mu=IS/c=evr/2c$, а механический момент $L=\mel vr$. Поскольку заряд электрона отрицателен, векторы
направлены противоположно (вывод~\ref{der:12:orb}):
\begin{equation}\label{eq:12:orb}
  \vect\mu_l=-\frac{e}{2\mel c}\vect L,\qquad \gamma_l=-\frac{e}{2\mel c},\qquad g_l=1 .
\end{equation}
Соотношение не зависит от формы орбиты и скорости — оно выполняется для любого движения заряда при
постоянном отношении $e/m$. В квантовой механике~\eqref{eq:12:orb} остаётся верным как операторное
равенство $\op{\vect\mu}_l=-\frac{e}{2\mel c}\op{\vect L}$. Следовательно,
\begin{equation}\label{eq:12:muz}
  \mu_{lz}=-\frac{e}{2\mel c}m\hbar=-\muB m,\quad m=-l,\dots,l;\qquad
  |\vect\mu_l|=\muB\sqrt{l(l+1)} .
\end{equation}
Проекция орбитального магнитного момента принимает $2l+1$ значений, кратных $\muB$.

\subsection{Атом в магнитном поле}

\paragraph{Энергия.} В однородном поле $\vect B\parallel z$ энергия $U=-\vect\mu\cdot\vect B=-\mu_zB=\muB Bm$. Уровень
с данным $l$ расщепляется на $2l+1$ равноотстоящих подуровней с интервалом $\muB B$ (\emph{нормальный эффект
Зеемана}). При $B=1$~Тл интервал $\approx5{,}8\cdot10^{-5}$~эВ — много меньше расстояний между уровнями атома.

\paragraph{Ларморова прецессия.} Момент сил $\vect\mu\times\vect B$ вызывает прецессию $\vect L$ вокруг $\vect B$ с частотой
(вывод~\ref{der:12:larmor})
\begin{equation}\label{eq:12:larmor}
  \omega_L=\frac{eB}{2\mel c},\qquad \hbar\omega_L=\muB B .
\end{equation}
Квантовое расщепление уровней на $\hbar\omega_L$ — квантовый аналог классической прецессии.

\paragraph{Сила в неоднородном поле.} В однородном поле на диполь действует только момент сил, в
неоднородном — ещё и сила
\begin{equation}\label{eq:12:force}
  \vect F=\nabla(\vect\mu\cdot\vect B)\quad\Longrightarrow\quad F_z\approx\mu_z\pdiff{B_z}{z}
\end{equation}
(когда поле направлено в основном вдоль $z$, а поперечные компоненты $\vect\mu$ из-за быстрой ларморовой
прецессии усредняются до нуля). Сила пропорциональна проекции $\mu_z$ — это и использовано в опыте
Штерна--Герлаха.

\subsection{Опыт Штерна--Герлаха (1922)}

\paragraph{Схема.} Пучок атомов серебра, испарявшихся в печи ($T\approx1300$~К), коллимировался щелями
и проходил между полюсами магнита специальной формы (один полюс — клин, другой — с желобом),
создающего сильно неоднородное поле с градиентом $\partial B_z/\partial z$ поперёк пучка (рис.~\ref{fig:12:sg}).
За магнитом атомы осаждались на стеклянной пластинке.

\paragraph{Ожидания.} Если ориентации магнитного момента распределены непрерывно (классика), $\mu_z$
принимает любые значения от $-\mu$ до $+\mu$, и след пучка должен размыться в сплошную полосу. Если проекция
квантуется, пучок должен расщепиться на конечное число компонент — по числу значений $\mu_z$.

\paragraph{Результат.} Пучок расщепился на \emph{две} компоненты, симметрично смещённые на расстояния,
соответствующие $\mu_z=\pm\muB$ (вывод~\ref{der:12:sg}). Без поля — одна полоса.

\paragraph{Интерпретация.}
\begin{enumerate}
 \item Опыт прямо доказал \emph{пространственное квантование} — дискретность проекции момента.
 \item Две компоненты невозможно объяснить орбитальным моментом: для него число компонент $2l+1$ нечётно.
       Более того, в основном состоянии серебра (конфигурация $[\mathrm{Kr}]4d^{10}5s^1$) единственный
       внешний электрон находится в $s$-состоянии, $l=0$, и орбитального момента нет вовсе.
 \item Вывод: электрон обладает собственным моментом с $2s+1=2$, т.\,е. $s=\frac12$, и магнитным моментом с
       проекцией $\pm\muB$. Отсюда $g_s=\mu_z/(\muB m_s)=2$. (В 1922~г. результат трактовали в рамках
       модели Бора; правильная интерпретация через спин появилась после 1925~г.)
\end{enumerate}
Аналогичные опыты с атомами водорода в основном состоянии (Фиппс и Тейлор, 1927; $l=0$) также дали две
компоненты, что исключает объяснение особенностями серебра.

\begin{figure}[htbp]
\centering
\begin{tikzpicture}[scale=0.92]
  % печь и щели
  \draw[thick, fill=gray!20] (-0.6,-0.5) rectangle (0.4,0.5); \node[below, font=\small] at (-0.1,-0.55) {печь};
  \draw[thick] (1.2,0.35) -- (1.2,0.07); \draw[thick] (1.2,-0.07) -- (1.2,-0.35);
  \draw[thick] (2.0,0.35) -- (2.0,0.05); \draw[thick] (2.0,-0.05) -- (2.0,-0.35);
  \node[below, font=\small] at (1.6,-0.4) {щели};
  % магнит
  \fill[accent!25] (3.0,0.35) -- (3.5,1.4) -- (6.5,1.4) -- (6.5,0.35) -- (6.1,0.35)-- (5.0,0.35) -- cycle;
  \draw[thick] (3.0,0.35) -- (3.5,1.4) -- (6.5,1.4);
  \node[font=\small] at (4.9,1.05) {N};
  \fill[keycol!25] (3.0,-1.4) rectangle (6.5,-0.55);
  \fill[white] (4.3,-0.85) rectangle (5.5,-0.55);
  \node[font=\small] at (4.0,-1.12) {S};
  \node[font=\small] at (4.9,-1.7) {магнит, $\partial B_z/\partial z\neq0$};
  % пучок
  \draw[very thick, gray] (0.4,0) -- (3.0,0);
  \draw[very thick, accent, ->] (3.0,0) .. controls (5.0,0) and (6.4,0.25) .. (8.6,0.85);
  \draw[very thick, accent, ->] (3.0,0) .. controls (5.0,0) and (6.4,-0.25) .. (8.6,-0.85);
  \draw[thin, dashed, gray] (6.5,0) -- (8.7,0);
  % экран
  \draw[very thick] (8.8,-1.5) -- (8.8,1.5); \node[above, font=\small] at (8.8,1.5) {пластинка};
  \node[right, font=\small, accent] at (8.9,0.85) {$\mu_z=+\muB$};
  \node[right, font=\small, accent] at (8.9,-0.85) {$\mu_z=-\muB$};
  \draw[->] (10.9,-0.5) -- (10.9,0.5) node[above, font=\small] {$z$};
  % ожидаемая картина
  \begin{scope}[xshift=12.4cm]
    \node[font=\scriptsize, align=center] at (0.3,1.6) {класс.};
    \fill[gray!45] (0.0,-1.0) rectangle (0.6,1.0);
    \node[font=\scriptsize, align=center] at (1.3,1.6) {опыт};
    \fill[accent] (1.0,0.7) rectangle (1.6,1.0);
    \fill[accent] (1.0,-1.0) rectangle (1.6,-0.7);
  \end{scope}
\end{tikzpicture}
\caption{Схема опыта Штерна--Герлаха. Справа: след пучка при классически непрерывной ориентации
моментов (сплошная полоса) и наблюдаемая картина (две полосы).}
\label{fig:12:sg}
\end{figure}

\subsection{Спин электрона}

Гипотезу спина выдвинули С.~Гаудсмит и Дж.~Уленбек (1925) для объяснения дублетной структуры
спектров щелочных металлов (жёлтая линия натрия — дублет $589{,}0$ и $589{,}6$~нм) и аномального эффекта
Зеемана. Свойства спина электрона:
\begin{align}
  &S^2=\hbar^2s(s+1)=\tfrac34\hbar^2,\qquad S_z=m_s\hbar=\pm\tfrac12\hbar, \label{eq:12:S}\\
  &\vect\mu_s=-g_s\frac{e}{2\mel c}\vect S=-g_s\muB\frac{\vect S}{\hbar},\qquad \mu_{sz}=-g_s\muB m_s\approx\mp\muB,\qquad g_s\approx2 . \label{eq:12:mus}
\end{align}
\emph{Гиромагнитное отношение для спина вдвое больше орбитального}: $\gamma_s=-g_se/2\mel c\approx-e/\mel c$.
Это было подтверждено опытами Эйнштейна--де Гааза (поворот намагничиваемого железного стержня),
давшими для ферромагнетиков $g\approx2$: ферромагнетизм обусловлен спинами, а не орбитальным движением.

Значение $g_s=2$ следует из релятивистского уравнения Дирака (1928). Квантовая электродинамика даёт
поправку: $g_s=2(1+\frac{\alpha}{2\pi}+\dots)=2{,}002\,319\dots$ — одно из наиболее точно проверенных
предсказаний физики.

\paragraph{Описание спина.} Волновая функция электрона зависит не только от $\vect r$, но и от спиновой
переменной $\sigma=\pm\frac12$: $\Psi(\vect r,\sigma)$, или двухкомпонентный столбец (спинор)
$\binom{\psi_+(\vect r)}{\psi_-(\vect r)}$, где $|\psi_\pm|^2$ — плотности вероятности для $S_z=\pm\hbar/2$. Операторы
проекций спина: $\op S_i=\frac\hbar2\sigma_i$ с матрицами Паули
\[
  \sigma_x=\begin{pmatrix}0&1\\1&0\end{pmatrix},\quad
  \sigma_y=\begin{pmatrix}0&-\ii\\\ii&0\end{pmatrix},\quad
  \sigma_z=\begin{pmatrix}1&0\\0&-1\end{pmatrix}.
\]

\paragraph{Полный момент и фактор Ланде.} Полный момент электрона $\vect J=\vect L+\vect S$, $j=l\pm\frac12$
(при $l=0$ только $j=\frac12$). Из-за различия $g_l$ и $g_s$ магнитный момент не коллинеарен $\vect J$; его
проекция на $\vect J$ определяется фактором Ланде (без вывода)
\begin{equation}\label{eq:12:lande}
  g_J=1+\frac{J(J+1)+S(S+1)-L(L+1)}{2J(J+1)},\qquad \mu_{Jz}=-g_J\muB M_J .
\end{equation}
Проверка: при $S=0$ $g_J=1$; при $L=0$ $g_J=2$. Расщепление уровня в слабом поле $\Delta E=g_J\muB BM_J$ —
\emph{аномальный эффект Зеемана}. Спин-орбитальное взаимодействие (энергия магнитного момента спина в
магнитном поле, возникающем в системе покоя электрона из-за его движения относительно ядра) расщепляет
уровни с $l\neq0$ на два ($j=l\pm\frac12$) — тонкая структура, например дублет натрия.

\section{Математические выводы}

\begin{derivation}[label=der:12:orb]{орбитальный магнитный момент}
\emph{Допущения:} электрон (заряд $-e$, масса $\mel$) движется равномерно по окружности радиуса $r$ со
скоростью $v$; гауссова система.

Сила тока — заряд, проходящий через сечение орбиты в единицу времени: $I=e/T=ev/(2\pi r)$.
Магнитный момент плоского витка: $\mu=IS/c=\dfrac{ev}{2\pi r}\cdot\dfrac{\pi r^2}{c}=\dfrac{evr}{2c}$.
Механический момент: $L=\mel vr$. Отношение
\[
  \frac{\mu}{L}=\frac{evr/2c}{\mel vr}=\frac{e}{2\mel c}.
\]
Ток электрона направлен против его скорости, поэтому $\vect\mu$ антипараллелен $\vect L$:
$\vect\mu_l=-\dfrac{e}{2\mel c}\vect L$. Для $L=\hbar$ получаем единицу магнитного момента — магнетон Бора
$\muB=e\hbar/2\mel c$.

\emph{Общность:} для произвольного движения $\vect\mu=\frac{1}{2c}\sum q\,\vect r\times\vect v$, а $\vect L=\sum m\,\vect r\times\vect v$;
при одинаковом $q/m$ для всех частиц $\vect\mu=\frac{q}{2mc}\vect L$.

\emph{Численно:} $\muB=\dfrac{4{,}803\cdot10^{-10}\cdot1{,}0546\cdot10^{-27}}{2\cdot9{,}109\cdot10^{-28}\cdot2{,}998\cdot10^{10}}=9{,}27\cdot10^{-21}$~эрг/Гс.
\end{derivation}

\begin{derivation}[label=der:12:larmor]{ларморова частота}
Уравнение моментов: $\dfrac{\dd\vect L}{\dd t}=\vect\mu\times\vect B=-\dfrac{e}{2\mel c}\vect L\times\vect B=\vect\omega_L\times\vect L$,
где $\vect\omega_L=\dfrac{e\vect B}{2\mel c}$. Уравнение вида $\dot{\vect L}=\vect\omega\times\vect L$ описывает вращение вектора $\vect L$ с
угловой скоростью $\vect\omega$ без изменения его длины и проекции на $\vect\omega$. Следовательно, $\vect L$ прецессирует
вокруг $\vect B$ с частотой $\omega_L=eB/2\mel c$. Энергия $\hbar\omega_L=\frac{e\hbar}{2\mel c}B=\muB B$ совпадает с
расстоянием между зеемановскими подуровнями $\Delta E=\muB B\,\Delta m$ при $\Delta m=1$.
\end{derivation}

\begin{derivation}[label=der:12:sg]{отклонение атомов в опыте Штерна--Герлаха}
\emph{Допущения:} атом движется вдоль $x$ со скоростью $v$; на участке длины $l_1$ внутри магнита действует
постоянная сила $F_z=\mu_z\,\partial B_z/\partial z$; далее — свободный пролёт длины $l_2$ до пластинки.

Внутри магнита (время $t_1=l_1/v$): ускорение $a=F_z/M$,
\[
  z_1=\frac{a t_1^2}{2}=\frac{\mu_z}{2M}\pdiff{B_z}{z}\frac{l_1^2}{v^2},\qquad v_z=at_1 .
\]
На пролёте: $z_2=v_z\frac{l_2}{v}=\frac{\mu_z}{M}\pdiff{B_z}{z}\frac{l_1l_2}{v^2}$. Полное смещение
\[
  z=\frac{\mu_z}{Mv^2}\pdiff{B_z}{z}\Bigl(\frac{l_1^2}{2}+l_1l_2\Bigr).
\]
Смещение пропорционально $\mu_z$: каждой проекции соответствует своя полоса. Для теплового пучка $Mv^2\sim3k_{\mathrm B}T$, поэтому смещение обратно пропорционально температуре печи; разброс скоростей в пучке размывает полосы, но не сливает их.

\emph{Оценка} (только магнит, $l_2=0$): $\partial B/\partial z\approx10^5$~Гс/см ($10^3$~Тл/м), $l_1=3{,}5$~см,
$v\approx5\cdot10^4$~см/с, $M=108\cdot1{,}66\cdot10^{-24}$~г$=1{,}8\cdot10^{-22}$~г, $\mu_z=\muB$:
\[
  z=\frac{9{,}27\cdot10^{-21}\cdot10^5\cdot(3{,}5)^2}{2\cdot1{,}8\cdot10^{-22}\cdot(5\cdot10^4)^2}\ \text{см}
  \approx1{,}3\cdot10^{-2}\ \text{см}=0{,}13\ \text{мм}.
\]
Расстояние между полосами $\sim0{,}2$--$0{,}3$~мм — порядок величины, наблюдавшийся в опыте. Из измеренного
смещения и известных $\partial B/\partial z$, $T$ найдено $|\mu_z|\approx\muB$.
\end{derivation}

\begin{derivation}[label=der:12:gs]{спиновый $g$-фактор из опыта}
Для атома серебра $L=0$, поэтому $\vect\mu=\vect\mu_s$. Наблюдаемые проекции $\mu_z=\pm\muB$, а $S_z=m_s\hbar$, $m_s=\pm\frac12$.
Из $\mu_z=-g_s\muB m_s$: $\muB=g_s\muB\cdot\frac12$ $\Rightarrow$ $g_s=2$. Гиромагнитное отношение
$|\mu_z/S_z|=\muB/(\hbar/2)=e/\mel c$ — вдвое больше орбитального.
\end{derivation}

\section{Итоговые формулы}

\keyformula{\vect\mu_l=-\frac{e}{2\mel c}\vect L,\qquad \mu_{lz}=-\muB m,\qquad \muB=\frac{e\hbar}{2\mel c}}
 {$\vect L$ — орбитальный момент, $m$ — магнитное квантовое число; $g_l=1$.}
 {орбитальное движение электрона (нерелятивистское).}
 {магнитный момент атома, зеемановское расщепление, отклонение в неоднородном поле.}

\keyformula{\vect\mu_s=-g_s\muB\frac{\vect S}{\hbar},\quad g_s\approx2{,}0023,\quad S_z=\pm\frac\hbar2,\quad \mu_{sz}\approx\mp\muB}
 {$\vect S$ — спин, $s=\frac12$.}
 {электрон; $g_s=2$ — по Дираку, поправка $\alpha/\pi$ — КЭД.}
 {спиновый магнетизм, опыт Штерна--Герлаха, ферромагнетизм.}

\keyformula{U=-\vect\mu\cdot\vect B,\qquad F_z=\mu_z\pdiff{B_z}{z},\qquad \omega_L=\frac{eB}{2\mel c}}
 {$\vect B$ — индукция поля, $\omega_L$ — ларморова частота.}
 {сила — в неоднородном поле; поперечные компоненты $\vect\mu$ усреднены прецессией.}
 {энергия в поле, расщепление пучков, прецессия момента.}

\keyformula{g_J=1+\frac{J(J+1)+S(S+1)-L(L+1)}{2J(J+1)},\qquad \Delta E=g_J\muB BM_J}
 {$L,S,J$ — орбитальное, спиновое и полное квантовые числа; $M_J$ — проекция $\vect J$.}
 {$LS$-связь, слабое поле (без вывода).}
 {аномальный эффект Зеемана.}

\section{Физическая интерпретация}

\textbf{Масштабы.} Магнитная энергия $\muB B\approx6\cdot10^{-5}$~эВ при 1~Тл мала по сравнению с
атомными энергиями (эВ) и тепловой энергией при комнатной температуре ($0{,}026$~эВ) — поэтому
парамагнетизм слаб, а ориентационная поляризация моментов требует низких температур или сильных полей.

\textbf{Ядерные моменты.} Для протона вместо $\mel$ стоит $m_p$: ядерный магнетон $\mu_N=e\hbar/2m_pc\approx\muB/1836$.
Магнитные моменты ядер малы, что объясняет малость сверхтонкого расщепления. Аномальные значения
моментов протона ($2{,}79\mu_N$) и нейтрона ($-1{,}91\mu_N$) указывают на их составную (кварковую) структуру.

\textbf{Значение опыта Штерна--Герлаха:} первое прямое доказательство квантования проекции момента и прототип
квантового измерения — каждый атом попадает лишь в одну из полос.

\section{Что необходимо знать на экзамене}

\begin{examcheck}
\checkhead{Определения}
\begin{itemize}
 \item магнитный момент, гиромагнитное отношение, $g$-фактор, магнетон Бора (формула и значение);
 \item спин: $s=\frac12$, $S^2=\frac34\hbar^2$, $S_z=\pm\frac\hbar2$; полный момент $\vect J=\vect L+\vect S$.
\end{itemize}
\checkhead{Формулы}
\begin{itemize}
 \item $\vect\mu_l=-\frac{e}{2\mel c}\vect L$, $g_l=1$; $\vect\mu_s=-g_s\muB\vect S/\hbar$, $g_s\approx2$;
 \item $U=-\vect\mu\vect B$, $F_z=\mu_z\partial B_z/\partial z$, $\omega_L=eB/2\mel c$; формула Ланде (для справки).
\end{itemize}
\checkhead{Выводы}
\begin{itemize}
 \item $\mu/L=e/2\mel c$ для круговой орбиты; ларморова прецессия;
 \item смещение атома в опыте Штерна--Герлаха; $g_s=2$ из наблюдения $\mu_z=\pm\muB$ при $L=0$.
\end{itemize}
\checkhead{Интерпретации}
\begin{itemize}
 \item почему две компоненты доказывают существование спина ($2s+1$ чётно, $l=0$ у Ag и H); опыт Эйнштейна--де Гааза.
\end{itemize}
\end{examcheck}
